\documentclass[reqno]{amsart}
\usepackage[T1]{fontenc}
\usepackage{amsmath,amssymb,amsthm}
\usepackage{xr-hyper}
\usepackage[hidelinks]{hyperref}
\externaldocument[ec-]{20261124T223000--elliptic-curves-for-the-proof__workbook}
\externaldocument[st-]{20261124T214500--supersingular-theta-divisibility__workbook}
\externaldocument[od-]{20261125T008000--one-cm-determinant__workbook}
\externaldocument[ip-]{20261124T214600--the-isotropic-plane__workbook}

% Chapter number; \ref{chap} gives it.
\newcounter{chap}
% Chapter numbers, generated from 20261125T012000--grh-workbook__guide.org
% by scripts/make-chapters.py.  Do not edit: change the order of the
% index and rerun the script (./preview-site.sh and CI do this).
%
% Usage in a sheet, after \newcounter{chap}:
%   % Chapter numbers, generated from 20261125T012000--grh-workbook__guide.org
% by scripts/make-chapters.py.  Do not edit: change the order of the
% index and rerun the script (./preview-site.sh and CI do this).
%
% Usage in a sheet, after \newcounter{chap}:
%   % Chapter numbers, generated from 20261125T012000--grh-workbook__guide.org
% by scripts/make-chapters.py.  Do not edit: change the order of the
% index and rerun the script (./preview-site.sh and CI do this).
%
% Usage in a sheet, after \newcounter{chap}:
%   \input{chapters}\chapnumber{<sheet ID>}
\makeatletter
\providecommand\chapters@set[2]{\@namedef{chapters@#1}{#2}}
\chapters@set{20261125T001000}{1}% zeros-to-ramp-sums
\chapters@set{20261125T001100}{2}% inert-primes-and-estimate-b
\chapters@set{20261125T005000}{3}% galois-fourier-and-averaging
\chapters@set{20261125T005500}{4}% columns-from-a-sphere
\chapters@set{20261125T000000}{5}% slope-method
\chapters@set{20261126T030100}{6}% a-twist-for-each-prime
\chapters@set{20261125T003000}{7}% the-chain-and-the-square-root
\chapters@set{20261125T010000}{8}% rank-lemma-and-zero-estimate
\chapters@set{20261125T010500}{9}% poincare-series-and-near-orthogonality
\chapters@set{20261124T223000}{10}% elliptic-curves-for-the-proof
\chapters@set{20261125T009500}{11}% fock-model-and-theta-intertwiners
\chapters@set{20261124T223100}{12}% complex-multiplication-and-inert-reduction
\chapters@set{20261125T008000}{13}% one-cm-determinant
\chapters@set{20261124T214500}{14}% supersingular-theta-divisibility
\chapters@set{20261124T214600}{15}% the-isotropic-plane
\chapters@set{20261125T013700}{16}% pairings-on-a-formal-group-of-height-two
\chapters@set{20261125T013500}{17}% hecke-nondegeneracy-on-the-tree
\chapters@set{20261126T030200}{18}% the-bound-is-attained
\chapters@set{20261125T009000}{19}% kronecker-limit-and-chowla-selberg
\chapters@set{20261125T002000}{20}% definite-theta-lifts
\chapters@set{20261125T015500}{21}% siegel-weil-in-the-smallest-case
\chapters@set{20261125T002100}{22}% definite-theta-lifts-and-norms
\chapters@set{20261125T016500}{23}% hecke-operators-and-theta-lifts
\chapters@set{20261125T016700}{24}% divisibility-along-the-chain
\chapters@set{20261125T004000}{25}% denominators-and-sha
\chapters@set{20261125T004100}{26}% tate-twists-and-sha
\chapters@set{20261126T020430}{27}% the-minus-part-by-gauss-sums
\chapters@set{20261125T014100}{28}% congruence-modules
\chapters@set{20261125T014500}{29}% eisenstein-congruences-and-selmer-classes
\chapters@set{20261125T014700}{30}% pseudocharacters-and-the-length-bound
\chapters@set{20261125T011000}{31}% galois-representations-black-box
\chapters@set{20261126T030300}{32}% the-data
\chapters@set{20261126T020210}{33}% bessel-poincare-series-and-their-gram-operators
\chapters@set{20261126T030400}{34}% the-seesaw-identity
\chapters@set{20261126T030500}{35}% divisibility-of-the-composite-lift
\chapters@set{20261126T030600}{36}% theorem-a-short-forms
\chapters@set{20261126T030700}{37}% the-determinant-of-a-chain
\chapters@set{20261126T030800}{38}% the-cost-of-one-lift
\chapters@set{20261126T030900}{39}% denominators-of-projections-to-lifts
\chapters@set{20261126T031000}{40}% twists-with-few-congruences
\chapters@set{20261126T031100}{41}% theorem-c-and-the-main-inequality
\providecommand\chapnumber[1]{%
  \@ifundefined{chapters@#1}%
    {\PackageError{chapters}{Sheet #1 is not listed in the index}%
      {List it in 20261125T012000--grh-workbook__guide.org\MessageBreak
       and run scripts/make-chapters.py.}}%
    {\edef\thechap{\@nameuse{chapters@#1}}}}
\makeatother
\chapnumber{<sheet ID>}
\makeatletter
\providecommand\chapters@set[2]{\@namedef{chapters@#1}{#2}}
\chapters@set{20261125T001000}{1}% zeros-to-ramp-sums
\chapters@set{20261125T001100}{2}% inert-primes-and-estimate-b
\chapters@set{20261125T005000}{3}% galois-fourier-and-averaging
\chapters@set{20261125T005500}{4}% columns-from-a-sphere
\chapters@set{20261125T000000}{5}% slope-method
\chapters@set{20261126T030100}{6}% a-twist-for-each-prime
\chapters@set{20261125T003000}{7}% the-chain-and-the-square-root
\chapters@set{20261125T010000}{8}% rank-lemma-and-zero-estimate
\chapters@set{20261125T010500}{9}% poincare-series-and-near-orthogonality
\chapters@set{20261124T223000}{10}% elliptic-curves-for-the-proof
\chapters@set{20261125T009500}{11}% fock-model-and-theta-intertwiners
\chapters@set{20261124T223100}{12}% complex-multiplication-and-inert-reduction
\chapters@set{20261125T008000}{13}% one-cm-determinant
\chapters@set{20261124T214500}{14}% supersingular-theta-divisibility
\chapters@set{20261124T214600}{15}% the-isotropic-plane
\chapters@set{20261125T013700}{16}% pairings-on-a-formal-group-of-height-two
\chapters@set{20261125T013500}{17}% hecke-nondegeneracy-on-the-tree
\chapters@set{20261126T030200}{18}% the-bound-is-attained
\chapters@set{20261125T009000}{19}% kronecker-limit-and-chowla-selberg
\chapters@set{20261125T002000}{20}% definite-theta-lifts
\chapters@set{20261125T015500}{21}% siegel-weil-in-the-smallest-case
\chapters@set{20261125T002100}{22}% definite-theta-lifts-and-norms
\chapters@set{20261125T016500}{23}% hecke-operators-and-theta-lifts
\chapters@set{20261125T016700}{24}% divisibility-along-the-chain
\chapters@set{20261125T004000}{25}% denominators-and-sha
\chapters@set{20261125T004100}{26}% tate-twists-and-sha
\chapters@set{20261126T020430}{27}% the-minus-part-by-gauss-sums
\chapters@set{20261125T014100}{28}% congruence-modules
\chapters@set{20261125T014500}{29}% eisenstein-congruences-and-selmer-classes
\chapters@set{20261125T014700}{30}% pseudocharacters-and-the-length-bound
\chapters@set{20261125T011000}{31}% galois-representations-black-box
\chapters@set{20261126T030300}{32}% the-data
\chapters@set{20261126T020210}{33}% bessel-poincare-series-and-their-gram-operators
\chapters@set{20261126T030400}{34}% the-seesaw-identity
\chapters@set{20261126T030500}{35}% divisibility-of-the-composite-lift
\chapters@set{20261126T030600}{36}% theorem-a-short-forms
\chapters@set{20261126T030700}{37}% the-determinant-of-a-chain
\chapters@set{20261126T030800}{38}% the-cost-of-one-lift
\chapters@set{20261126T030900}{39}% denominators-of-projections-to-lifts
\chapters@set{20261126T031000}{40}% twists-with-few-congruences
\chapters@set{20261126T031100}{41}% theorem-c-and-the-main-inequality
\providecommand\chapnumber[1]{%
  \@ifundefined{chapters@#1}%
    {\PackageError{chapters}{Sheet #1 is not listed in the index}%
      {List it in 20261125T012000--grh-workbook__guide.org\MessageBreak
       and run scripts/make-chapters.py.}}%
    {\edef\thechap{\@nameuse{chapters@#1}}}}
\makeatother
\chapnumber{<sheet ID>}
\makeatletter
\providecommand\chapters@set[2]{\@namedef{chapters@#1}{#2}}
\chapters@set{20261125T001000}{1}% zeros-to-ramp-sums
\chapters@set{20261125T001100}{2}% inert-primes-and-estimate-b
\chapters@set{20261125T005000}{3}% galois-fourier-and-averaging
\chapters@set{20261125T005500}{4}% columns-from-a-sphere
\chapters@set{20261125T000000}{5}% slope-method
\chapters@set{20261126T030100}{6}% a-twist-for-each-prime
\chapters@set{20261125T003000}{7}% the-chain-and-the-square-root
\chapters@set{20261125T010000}{8}% rank-lemma-and-zero-estimate
\chapters@set{20261125T010500}{9}% poincare-series-and-near-orthogonality
\chapters@set{20261124T223000}{10}% elliptic-curves-for-the-proof
\chapters@set{20261125T009500}{11}% fock-model-and-theta-intertwiners
\chapters@set{20261124T223100}{12}% complex-multiplication-and-inert-reduction
\chapters@set{20261125T008000}{13}% one-cm-determinant
\chapters@set{20261124T214500}{14}% supersingular-theta-divisibility
\chapters@set{20261124T214600}{15}% the-isotropic-plane
\chapters@set{20261125T013700}{16}% pairings-on-a-formal-group-of-height-two
\chapters@set{20261125T013500}{17}% hecke-nondegeneracy-on-the-tree
\chapters@set{20261126T030200}{18}% the-bound-is-attained
\chapters@set{20261125T009000}{19}% kronecker-limit-and-chowla-selberg
\chapters@set{20261125T002000}{20}% definite-theta-lifts
\chapters@set{20261125T015500}{21}% siegel-weil-in-the-smallest-case
\chapters@set{20261125T002100}{22}% definite-theta-lifts-and-norms
\chapters@set{20261125T016500}{23}% hecke-operators-and-theta-lifts
\chapters@set{20261125T016700}{24}% divisibility-along-the-chain
\chapters@set{20261125T004000}{25}% denominators-and-sha
\chapters@set{20261125T004100}{26}% tate-twists-and-sha
\chapters@set{20261126T020430}{27}% the-minus-part-by-gauss-sums
\chapters@set{20261125T014100}{28}% congruence-modules
\chapters@set{20261125T014500}{29}% eisenstein-congruences-and-selmer-classes
\chapters@set{20261125T014700}{30}% pseudocharacters-and-the-length-bound
\chapters@set{20261125T011000}{31}% galois-representations-black-box
\chapters@set{20261126T030300}{32}% the-data
\chapters@set{20261126T020210}{33}% bessel-poincare-series-and-their-gram-operators
\chapters@set{20261126T030400}{34}% the-seesaw-identity
\chapters@set{20261126T030500}{35}% divisibility-of-the-composite-lift
\chapters@set{20261126T030600}{36}% theorem-a-short-forms
\chapters@set{20261126T030700}{37}% the-determinant-of-a-chain
\chapters@set{20261126T030800}{38}% the-cost-of-one-lift
\chapters@set{20261126T030900}{39}% denominators-of-projections-to-lifts
\chapters@set{20261126T031000}{40}% twists-with-few-congruences
\chapters@set{20261126T031100}{41}% theorem-c-and-the-main-inequality
\providecommand\chapnumber[1]{%
  \@ifundefined{chapters@#1}%
    {\PackageError{chapters}{Sheet #1 is not listed in the index}%
      {List it in 20261125T012000--grh-workbook__guide.org\MessageBreak
       and run scripts/make-chapters.py.}}%
    {\edef\thechap{\@nameuse{chapters@#1}}}}
\makeatother
\chapnumber{20261124T223100}
\renewcommand{\thesection}{\thechap.\arabic{section}}

\theoremstyle{definition}
\newtheorem{problem}{Problem}[section]
\theoremstyle{plain}
\newtheorem{theorem}{Theorem}[section]

\newcommand{\Z}{\mathbb{Z}}
\newcommand{\Q}{\mathbb{Q}}
\newcommand{\C}{\mathbb{C}}
\newcommand{\F}{\mathbb{F}}
\newcommand{\ii}{\mathrm{i}}
\newcommand{\Oc}{\mathcal{O}}
\newcommand{\End}{\operatorname{End}}

\setlength{\emergencystretch}{2em}

\title{Complex multiplication and reduction at an inert prime}
\author{}
\date{November 24, 2026}

\begin{document}
\maketitle
\refstepcounter{chap}\label{chap}

\paragraph{Prerequisites.}
Chapter~\ref{ec-chap}, \emph{Complex tori, theta functions and periods}, Section~\ref{ec-sec:tori}: the Weierstrass functions, Problems~\ref{ec-pr:zeta-add} and~\ref{ec-pr:square}, and the lemniscatic lattice \(\varpi\Z[\ii]\) of the curve \(E \colon y^2 = x^3 - x\).  Orders in imaginary quadratic fields.  Elliptic curves over a field at the level of \cite[Ch.~III]{Si1}, and the statements, not the proofs, of \cite[Ch.~IV, \S\S V.1--V.4, Ch.~VI]{Si1}.

\section{Complex multiplication}\label{sec:cm}

Which tori have endomorphisms other than \(u \mapsto nu\)?  Every holomorphic map \(\C/\Lambda \to \C/\Lambda\) that fixes \(0\) is \(u \mapsto \alpha u\) with \(\alpha\Lambda \subseteq \Lambda\) \cite[Thm.~VI.4.1]{Si1}.  So if \(E\) corresponds to \(\Lambda\), then \(\End(E) = \{\alpha \in \C : \alpha\Lambda \subseteq \Lambda\}\), and we write \([\alpha]\) for \(u \mapsto \alpha u\).

\begin{problem}\label{pr:cm-order}
If \(\End(E) \ne \Z\), then \(K := \Q(\alpha)\) is imaginary quadratic for every \(\alpha \in \End(E) \smallsetminus \Z\), \(\End(E)\) is an order in \(K\), and \(\Lambda\) is homothetic to a lattice contained in \(K\).
\end{problem}

We then say that \(E\) has \emph{complex multiplication} by that order.

\begin{problem}\label{pr:cm-examples}
\(\End(\C/\Z[\ii]) = \Z[\ii]\), \(\End(\C/\Z[2\ii]) = \Z[2\ii]\), and \(\End(\C/\Z[\rho]) = \Z[\rho]\) with \(\rho = e^{2\pi\ii/3}\).
\end{problem}

\begin{problem}\label{pr:deg}
For \(\alpha \in \End(E)\), \(\alpha \ne 0\), the kernel of \([\alpha]\) on \(\C/\Lambda\) has \(\lvert\alpha\rvert^2\) elements.
\end{problem}

For each nonzero \(\alpha\in\End(E)\), the function \(\wp(\alpha u)\) is even and elliptic for \(\Lambda\), so it is a rational function of \(\wp(u)\), and \([\alpha]\) is a rational map of the curve.

\begin{problem}\label{pr:one-plus-i}
For the lemniscatic lattice, \(\wp((1+\ii)u) = \dfrac{\wp(u)^2 - 1}{2\ii\, \wp(u)}\).
\end{problem}

\begin{problem}\label{pr:i-action}
On \(E \colon y^2 = x^3 - x\), \([\ii](x, y) = (-x, \ii y)\).  For every \(\alpha \in \End(E)\), \([\alpha]^*(dx/2y) = \alpha\, dx/2y\).
\end{problem}

\begin{problem}\label{pr:j1728}
For \(a, b \in \C^\times\), the curve \(y^2 = x^3 + ax\) has complex multiplication by \(\Z[\ii]\), and \(y^2 = x^3 + b\) has complex multiplication by \(\Z[\rho]\).
\end{problem}

The endomorphisms of \(E \colon y^2 = x^3 - x\) extend to its elliptic curve model over \(\Z[\ii,\tfrac12]\), so they can be reduced at each prime of this ring above an odd prime \(p\).  What do they look like near the origin?  In a local parameter at the origin, addition gives a formal group law, and each endomorphism acts on it by a power series.

\section{The formal group}\label{sec:formal}

Let \(E\) be given by a Weierstrass equation over a ring \(R\), and let \(z := -x/y\).  The group law, written in \(z\), is a formal group law \(F(z_1, z_2) \in R[[z_1, z_2]]\), and every endomorphism defined over \(R\) becomes a series \([\alpha](z) \in R[[z]]\) \cite[\S\S IV.1--IV.2]{Si1}.  The invariant differential becomes \(\omega(z) = \sum_{n \ge 1} c_n z^{n-1}\, d z\) with \(c_1 = 1\) and \(c_n \in R\).  If \(R\) is torsion-free, then over \(R \otimes_\Z \Q\) the \emph{logarithm} \(\ell(z) := \sum_{n \ge 1} c_n z^n/n\) satisfies \(F(z_1, z_2) = \ell^{-1}(\ell(z_1) + \ell(z_2))\) \cite[\S\S IV.4--IV.5]{Si1}.  Over \(\C\), \(\ell(z(u)) = u\).  The \emph{invariant derivation} is \(D := \ell'(z)^{-1}\, d/dz\), where \(\ell'(z) = \sum_{n \ge 1} c_n z^{n-1}\) has the coefficients of \(\omega\); over \(\C\), \(D = d/du\).

For a one-parameter formal group over \(\Z_p\), either \([p](z) \equiv 0 \pmod p\) or \([p](z) \equiv g(z^{p^h}) \pmod p\) with \(g'(0) \not\equiv 0\); the integer \(h\) is the \emph{height} \cite[\S IV.7]{Si1}.  For an elliptic curve with good reduction at \(p\), \(h \in \{1, 2\}\), and \(h = 2\) exactly when the reduction is supersingular \cite[Thm.~V.3.1]{Si1}.  The two extreme models are \(\hat G_a(X, Y) = X + Y\) and \(\hat G_m(X, Y) = X + Y + XY\).

\begin{problem}\label{pr:Gm}
For \(\hat G_m\), \(\ell(X) = \log(1+X)\), \([p](X) \equiv X^p \pmod p\), and the height is \(1\).
\end{problem}

\begin{problem}\label{pr:Dp-models}
As operators on \(\Z_p[[X]]\) modulo \(p\): for \(\hat G_a\), \(D^p \equiv 0\); for \(\hat G_m\), \(D = (1+X)\,d/dX\) and \(D^p \equiv D\).
\end{problem}

For a one-parameter formal group over \(\Z_p\), its normalized invariant derivation satisfies \(D^p \equiv \lambda D \pmod p\) for some \(\lambda \in \F_p\).  Problems~\ref{st-pr:Dp-derivation} and~\ref{st-pr:invariant} of Chapter~\ref{st-chap} give the argument for the lemniscate, and the same argument applies to every one-parameter formal group over \(\Z_p\).  Which constant is \(\lambda\)?

\begin{problem}\label{pr:lambda}
Let \(\hat G\) be the formal group of a Weierstrass equation over \(\Z_{(p)}\), with \(D^p \equiv \lambda D \pmod p\).  Then \((D^n z)(0) = n!\, [u^n]\, \ell^{-1}(u)\) for every \(n\), and \(\lambda \equiv c_p \pmod p\).
\end{problem}

\begin{problem}\label{pr:coef-p}
For \(\hat G\) as in Problem~\ref{pr:lambda}, the coefficient of \(z^p\) in \([p](z)\) is congruent to \(c_p\) modulo \(p\).  Hence the height is \(1\) exactly when \(p \nmid c_p\).
\end{problem}

Now we compute \(c_n\) for \(E \colon y^2 = x^3 - x\).  Let \(w := -1/y\).

\begin{problem}\label{pr:w-z}
For \(y^2 = x^3 - x\), \(w + z w^2 = z^3\), and \(\omega(z) = d z / \sqrt{1 + 4z^4}\).
\end{problem}

\begin{problem}\label{pr:log}
For \(y^2 = x^3 - x\), \(\ell(z) = \sum_{m \ge 0} (-1)^m \binom{2m}{m} z^{4m+1}/(4m+1)\).  For odd \(n\), \(c_n\) is the coefficient of \(x^{n-1}\) in \((x^3 - x)^{(n-1)/2}\).
\end{problem}

For odd \(p\) and a curve \(y^2 = f(x)\) over \(\F_p\) with \(f\) a cubic with distinct roots over \(\overline{\F}_p\), the coefficient \(A_p\) of \(x^{p-1}\) in \(f^{(p-1)/2}\) is the \emph{Hasse invariant}, and the curve is supersingular exactly when \(A_p = 0\) \cite[Thm.~V.4.1(a)]{Si1}.  By Problem~\ref{pr:log}, \(A_p\) is the reduction of \(c_p\) modulo \(p\) for \(y^2 = x^3 - x\), and Problems~\ref{pr:lambda} and~\ref{pr:coef-p} recover the criterion for this curve.  When \(p \mid c_p\), the first term of \([p](z)\) modulo \(p\) lies beyond \(z^p\).  What is its first nonzero term?

\section{Reduction at an inert prime}\label{sec:reduction}

Which primes are supersingular for a curve with complex multiplication?  Deuring's answer \cite{Deu41}: if \(E\) has complex multiplication by an order in \(K\) and good reduction at a prime above \(p\), the reduction is supersingular if and only if \(p\) does not split in \(K\).  For \(y^2 = x^3 - x\) these are the \(p \equiv 3 \pmod 4\).  We check this by hand.

\begin{problem}\label{pr:count}
Let \(p \equiv 3 \pmod 4\).  Then \(\#E(\F_p) = p + 1\), that is, \(a_p := p + 1 - \#E(\F_p) = 0\).
\end{problem}

\begin{problem}\label{pr:hasse-x3x}
For \(p = 4m+1\), \(A_p = (-1)^m \binom{2m}{m} \not\equiv 0 \pmod p\).  For \(p \equiv 3 \pmod 4\), \(A_p = 0\).
\end{problem}

For \(p \equiv 1 \pmod 4\) one has \(a_p = \gamma + \bar\gamma\), where \(\gamma \in \Z[\ii]\) has \(\gamma\bar\gamma = p\) and \(\gamma \equiv 1 \pmod{2 + 2\ii}\) \cite[Ch.~18]{IR}.

\begin{problem}\label{pr:ap-split}
\(a_5 = -2\) and \(a_{13} = 6\), in agreement with \(\gamma = -1 + 2\ii\) and \(\gamma = 3 + 2\ii\).
\end{problem}

Now let \(p \equiv 3 \pmod 4\), let \(\phi(x, y) := (x^p, y^p)\) be the Frobenius endomorphism of \(E\) over \(\F_p\), and reduce \([\ii]\) modulo \(p\) into \(\F_{p^2}\).

\begin{problem}\label{pr:frob}
\(\phi^2 = [-p]\) and \(\phi \circ [\ii] = [-\ii] \circ \phi\).  Hence \(\End(E_{\bar\F_p})\) is not commutative.
\end{problem}

\begin{problem}\label{pr:p-series}
Let \(p \equiv 3 \pmod 4\).  In the parameter \(z = -x/y\), \([p](z) \equiv -z^{p^2} \pmod p\), and \([\ii](z) = \ii z\).
\end{problem}

For instance \([3](z) \equiv -z^9 \pmod 3\).  To describe the formal group at these primes, we use a theorem of Lubin and Tate \cite{LT}.  Let \(\Oc\) be the ring of integers of a finite extension of \(\Q_p\), with uniformizer \(\pi\) and residue field of order \(q\).  Lubin and Tate prove:
(i) each \(f \in \Oc[[X]]\) with \(f(X) \equiv \pi X \bmod X^2\) and \(f(X) \equiv X^q \bmod \pi\) is an endomorphism of a unique formal group \(F_f\) over \(\Oc\);
(ii) each \(a \in \Oc\) gives a unique endomorphism \([a]_f\) of \(F_f\) with linear term \(aX\), and \(a \mapsto [a]_f\) is a ring homomorphism;
(iii) for two such \(f\), \(g\) there is a unique isomorphism \(\theta \colon F_f \to F_g\) over \(\Oc\) with linear term \(X\), and it commutes with the \([a]\).
By Hazewinkel's functional-equation lemma \cite[I.2.2]{Ha}, the \emph{model logarithm} \(\ell_\pi(X) := X + X^q/\pi + X^{q^2}/\pi^2 + \dotsb\) is the logarithm of such an \(F_f\).

Let \(p \equiv 3 \pmod 4\) and \(\Oc := \Z_p[\ii]\), so that \(\pi := -p\) is a uniformizer and \(q = p^2\).

\begin{problem}\label{pr:LT}
The series \(f := [-p]\) satisfies the hypotheses of (i), the formal group \(\hat E\) of \(y^2 = x^3 - x\) is \(F_f\), and \([a]_f\) extends \(\alpha \mapsto [\alpha]\) from \(\Z[\ii]\) to \(\Oc\).
\end{problem}

\begin{problem}\label{pr:model-log}
There is \(\theta \in z + z^2\Oc[[z]]\) with \(\ell = \ell_{-p} \circ \theta\).  Consequently the coefficients of \(\ell\) in degrees less than \(p^2\) are \(p\)-integral, and \(p \cdot [z^{p^2}]\, \ell \equiv -1 \pmod p\).
\end{problem}

For \(p = 3\): \([z^9]\,\ell = \tfrac23\), and \(3 \cdot \tfrac23 \equiv -1 \pmod 3\).  Chapter~\ref{st-chap} proves the same congruence for the lemniscate with binomial coefficients (Problem~\ref{st-pr:log-q}).

In the model logarithm the denominators occur only in the degrees \(q, q^2, \dotsc\), and there they are \(p, p^2, \dotsc\).  For the lemniscate, Theorem~\ref{st-thm:single} of Chapter~\ref{st-chap} gives a divisibility bound for the invariant derivatives of series integral in the formal parameter.  Problem~\ref{st-pr:failure} of that chapter gives a theta function whose fourth derivative at the origin is \(4\) when \(p=3\).  Chapter~\ref{od-chap} uses \(D^p\) in a single determinant; theta functions are the subject of Chapters~\ref{st-chap} and~\ref{ip-chap}.  This chapter and those three prepare the divisibility of derivatives of theta functions at primes that do not split in the imaginary quadratic field of the proof.  The corresponding statement of the rewritten Lean proof is \texttt{Parameters.norm\_jet\_minorProd\_jetSeries\_le}; the factorization of one theta series on which it rests is \texttt{CMLattice.exists\_exponent\_tsum}.

\section*{Hints}

\subsection*{\S\ref{sec:cm}}
\begin{description}
\item[\ref{pr:cm-order}] Write \(\alpha\omega_1 = a\omega_1 + b\omega_2\) and \(\alpha\omega_2 = c\omega_1 + d\omega_2\) with integers \(a, b, c, d\).  Then \(\alpha\) is a root of \(x^2 - (a+d)x + ad - bc\).  If \(\alpha \notin \mathbb{R}\), then \(\omega_2/\omega_1 \in \Q(\alpha)\); if \(\alpha \in \mathbb{R}\), then \(b = c = 0\).
\item[\ref{pr:cm-examples}] \(\alpha \cdot 1 \in \Lambda\).
\item[\ref{pr:deg}] Compare the areas of fundamental parallelograms for \(\Lambda\) and \(\alpha\Lambda\).
\item[\ref{pr:one-plus-i}] Apply Problem~\ref{ec-pr:zeta-add} with \(v = \ii u\), using Problem~\ref{ec-pr:square} and \(\wp'^2 = 4\wp^3 - 4\wp\).
\item[\ref{pr:i-action}] \(x = \wp(u)\) and \(y = \wp'(u)/2\).
\item[\ref{pr:j1728}] Find automorphisms of the curves of order \(4\) and \(3\).
\end{description}

\subsection*{\S\ref{sec:formal}}
\begin{description}
\item[\ref{pr:Gm}] \([p](X) = (1+X)^p - 1\).
\item[\ref{pr:Dp-models}] \(D(1+X)^k = k(1+X)^k\).  For \(\hat G_a\) and \(n\ge p\), \((d/dX)^p X^n = n(n-1)\dotsm(n-p+1)X^{n-p}\).  For \(0\le n<p\), the derivative is zero.
\item[\ref{pr:lambda}] Use the coordinate \(u=\ell(z)\) for the first assertion.  For the second, apply the operator congruence to the truncation of \(\ell\) in degrees less than \(p\).
\item[\ref{pr:coef-p}] Compare \(\ell([p](z)) = p\,\ell(z)\) modulo \((p, z^{p+1})\).
\item[\ref{pr:w-z}] Substitute \(x = z/w\) and \(y = -1/w\).  Solve the quadratic for \(w\), and compute \(dx/2y = -\tfrac{z}{2}\, d\log x\).
\item[\ref{pr:log}] The binomial series of \((1 + 4z^4)^{-1/2}\).  For the second statement, \((x^3 - x)^{2m} = x^{2m}(x^2 - 1)^{2m}\).
\end{description}

\subsection*{\S\ref{sec:reduction}}
\begin{description}
\item[\ref{pr:count}] \(x^3 - x\) is odd, and \(-1\) is not a square modulo \(p\).
\item[\ref{pr:hasse-x3x}] Problem~\ref{pr:log}.  For the non-vanishing, \(2m < p\).
\item[\ref{pr:ap-split}] Count the solutions of \(y^2 = x^3 - x\) over \(\F_5\) and \(\F_{13}\).
\item[\ref{pr:frob}] \(\phi^2 - a_p\phi + p = 0\) in \(\End(E)\) \cite[\S V.2]{Si1}, and Problem~\ref{pr:count}.  For the second relation, \(\ii^p = -\ii\) in \(\F_{p^2}\).
\item[\ref{pr:p-series}] In the parameter \(z\), \(\phi\) acts on the formal group over \(\F_p\) by \(z \mapsto z^p\), and \([-1](z) = -z\).
\item[\ref{pr:LT}] Problem~\ref{pr:p-series}.  For the last statement, use the uniqueness in (ii).
\item[\ref{pr:model-log}] Use (iii).  A formal group over a \(\Q\)-algebra has only one logarithm with linear term \(z\).  For the consequence, expand \(\ell_{-p}(\theta(z))\) up to degree \(q\).
\end{description}

\begin{thebibliography}{Deu}
\bibitem[Deu]{Deu41} M.~Deuring, \emph{Die Typen der Multiplikatorenringe elliptischer Funktionenk\"orper}, Abh.\ Math.\ Sem.\ Hansischen Univ.\ \textbf{14} (1941), 197--272.
\bibitem[Ha]{Ha} M.~Hazewinkel, \emph{Formal Groups and Applications}, Academic Press, 1978.
\bibitem[IR]{IR} K.~Ireland and M.~Rosen, \emph{A Classical Introduction to Modern Number Theory}, 2nd ed., Grad.\ Texts in Math.\ 84, Springer, 1990.
\bibitem[LT]{LT} J.~Lubin and J.~Tate, \emph{Formal complex multiplication in local fields}, Ann.\ of Math.\ \textbf{81} (1965), 380--387.
\bibitem[Si1]{Si1} J.~H.~Silverman, \emph{The Arithmetic of Elliptic Curves}, 2nd ed., Grad.\ Texts in Math.\ 106, Springer, 2009.
\end{thebibliography}

\end{document}
